\clearpage
\begin{trapbox}
\textbf{\textcolor{CTrap}{坑一：看到 $4|a|$，就写成“最小正周期”}}
\begin{center}\begin{tikzpicture}[font=\small,>=stealth,line join=round]
\draw[gray!40,rounded corners=3pt] (0,0) rectangle (14.2,8.8);\node at (7.1,8.200000000000001) {例5：同样奇函数、同一条轴，最小周期仍可不同};
\node at (3.1500000000000004,7.3) {$p=\sin\frac{\pi x}2$};\begin{scope}[shift={(1.19,4.338709677419355)},x=0.9800000000000001cm,y=1.225806451612903cm]\draw[->,gray] (-0.5,0)--(4.5,0) node[right]{$x$};\draw[->,gray] (0,-1.5)--(0,1.6) node[above]{$y$};\draw[blue,thick,solid,domain=-0.3:4.3,samples=180] plot(\x,{sin(deg(pi*\x/2))});\fill[orange!85!black] (1,1) circle(2pt);\fill[orange!85!black] (3,-1) circle(2pt);\draw[violet,dashed] (1,-1.3)--(1,1.3);\end{scope}\node[align=center,text width=5.8cm] at (3.1500000000000004,1.7) {关于 $x=1$ 对称\\最小正周期4};\node at (9.95,7.3) {$q=\sin\frac{3\pi x}2$};\begin{scope}[shift={(7.99,4.338709677419355)},x=0.9800000000000001cm,y=1.225806451612903cm]\draw[->,gray] (-0.5,0)--(4.5,0) node[right]{$x$};\draw[->,gray] (0,-1.5)--(0,1.6) node[above]{$y$};\draw[violet,thick,solid,domain=-0.3:4.3,samples=180] plot(\x,{sin(deg(3*pi*\x/2))});\fill[orange!85!black] (1,-1) circle(2pt);\draw[violet,dashed] (1,-1.3)--(1,1.3);\end{scope}\node[align=center,text width=5.8cm] at (9.95,1.7) {也关于 $x=1$ 对称\\最小正周期 $\frac43$};
\node[text width=12.6cm,align=center] at (7.1,.65) {两图同尺度：相同条件保证4都是周期，却不能保证最小周期相同。};
\end{tikzpicture}\end{center}



\textbf{错在哪里：}由题设只证明了 $f(x+4a)=f(x)$，即 $4|a|$ 是\textbf{一个}周期，不能排除更短的周期。比如 $\sin\frac{3\pi x}{2}$ 既是奇函数又关于 $x=1$ 对称，但最小正周期是 $4/3$，而不是 $4$。

\textbf{应该怎样写：}“$4|a|$ 是一个正周期”；只有额外证明不存在更小的正周期，才写“最小”。

\medskip

\end{trapbox}
\clearpage
\begin{trapbox}
\textbf{\textcolor{CTrap}{坑二：认为 $2|a|$ 就是周期}}
\begin{center}\begin{tikzpicture}[font=\small,>=stealth,line join=round]
\draw[gray!40,rounded corners=3pt] (0,0) rectangle (14.2,8.3);\node at (7.1,7.700000000000001) {易错二：5先变负5，再回到5；隔2尚未还原};
\begin{scope}[shift={(1.836842105263158,3.8706896551724133)},x=1.0526315789473684cm,y=0.3172413793103448cm]\draw[->,gray] (-0.7,0)--(5,0) node[right]{$x$};\draw[->,gray] (0,-7)--(0,7.5) node[above]{$y$};\draw[teal,thick,solid,domain=-0.3:4.6,samples=180] plot(\x,{5*cos(deg(pi*\x/2))});\fill[orange!85!black] (0,5) circle(2.2pt);\fill[orange!85!black] (2,-5) circle(2.2pt);\fill[orange!85!black] (4,5) circle(2.2pt);\end{scope}\node[anchor=west,align=left,text width=5.0cm] at (8.4,4.0) {示例 $f=5\cos\frac{\pi x}2$\\[9pt]偶函数，中心 $(1,0)$\\[9pt]$f(0)=5$\\[9pt]$f(2)=-5\ne5$\\[9pt]$f(4)=5$};
\node[text width=12.6cm,align=center] at (7.1,.65) {本图直接使用误区中的数值5；反号不等于相等，非零函数不能把2当周期。};
\end{tikzpicture}\end{center}



本条专指奇配轴、偶配点。题目推出的是 $f(x+2a)=-f(x)$，右边多了一个负号！如果某处 $f(x)=5$，移动后就变成 $-5$，不能说整张图没变化。只要函数不是处处为零，$2|a|$ 就不是周期；再移动一次才恢复原值。

\medskip

\end{trapbox}
\clearpage
\begin{trapbox}
\textbf{\textcolor{CTrap}{坑三：把“奇配轴、偶配点”和另外两种混为一谈}}
\begin{center}\begin{tikzpicture}[font=\small,>=stealth,line join=round]
\draw[gray!40,rounded corners=3pt] (0,0) rectangle (14.2,8.8);\node at (7.1,8.200000000000001) {同类对称只走2就还原，与奇配轴、偶配点不同};
\node at (3.1500000000000004,7.3) {奇配点：$f=\sin\pi x$};\begin{scope}[shift={(1.4795454545454545,4.338709677419355)},x=1.1136363636363635cm,y=1.225806451612903cm]\draw[->,gray] (-0.7,0)--(3.7,0) node[right]{$x$};\draw[->,gray] (0,-1.5)--(0,1.6) node[above]{$y$};\draw[blue,thick,solid,domain=-0.4:3.4,samples=180] plot(\x,{sin(deg(pi*\x))});\fill[orange!85!black] (0,0) circle(2pt);\fill[orange!85!black] (1,0) circle(2pt);\end{scope}\node[align=center,text width=5.8cm] at (3.1500000000000004,1.7) {中心 $(0,0),(1,0)$\\一个周期为2};\node at (9.95,7.3) {偶配轴：$g=\cos\pi x$};\begin{scope}[shift={(8.279545454545454,4.338709677419355)},x=1.1136363636363635cm,y=1.225806451612903cm]\draw[->,gray] (-0.7,0)--(3.7,0) node[right]{$x$};\draw[->,gray] (0,-1.5)--(0,1.6) node[above]{$y$};\draw[teal,thick,solid,domain=-0.4:3.4,samples=180] plot(\x,{cos(deg(pi*\x))});\foreach \a in {0,1}{\draw[violet,dashed] (\a,-1.3)--(\a,1.3);}\end{scope}\node[align=center,text width=5.8cm] at (9.95,1.7) {轴 $x=0,1$\\一个周期也为2};
\node[text width=12.6cm,align=center] at (7.1,.65) {两图同尺度：同类组合两次都不变号或两次变号，移2就保持原高度。};
\end{tikzpicture}\end{center}



奇函数 $+$ 对称轴、偶函数 $+$ 对称中心：移动 $2a$ 后\textbf{变号}，周期为 $4|a|$。

奇函数 $+$ 对称中心、偶函数 $+$ 对称轴：移动 $2a$ 后\textbf{不变号}，因此 $2|a|$ 已是一个周期。两类都可能有更短的最小正周期，千万别把“一个周期”认作“最小周期”。

\medskip

\end{trapbox}
\clearpage
\begin{trapbox}
\textbf{\textcolor{CTrap}{坑四：$a$ 是负数，就得到负周期}}
\begin{center}\begin{tikzpicture}[font=\small,>=stealth,line join=round]
\draw[gray!40,rounded corners=3pt] (0,0) rectangle (14.2,8.3);\node at (7.1,7.700000000000001) {例6：轴在负1，向左走4与向右走4都能重复};
\begin{scope}[shift={(4.1,3.8842857142857143)},x=0.6cm,y=1.3142857142857143cm]\draw[->,gray] (-5,0)--(5,0) node[right]{$x$};\draw[->,gray] (0,-1.7)--(0,1.8) node[above]{$y$};\draw[blue,thick,solid,domain=-4.5:4.5,samples=180] plot(\x,{sin(deg(pi*\x/2))});\fill[orange!85!black] (-1,-1) circle(2.2pt);\fill[orange!85!black] (0,0) circle(2.2pt);\fill[orange!85!black] (2,0) circle(2.2pt);\fill[orange!85!black] (4,0) circle(2.2pt);\draw[violet,dashed] (-1,-1.4)--(-1,1.4);\end{scope}\node[anchor=west,align=left,text width=5.0cm] at (8.4,4.0) {$h=\sin\frac{\pi x}2$\\[9pt]奇函数，轴 $x=-1$\\[9pt]$h(x-2)=-h(x)$\\[9pt]$h(x+4)=h(x)$\\[9pt]正周期可取4};
\node[text width=12.6cm,align=center] at (7.1,.65) {a为负时，移动方向变了，正的重复距离仍为4|a|。};
\end{tikzpicture}\end{center}



$a$ 可能为负，但正周期一定是正数。由 $f(x+4a)=f(x)$，可以得到 $f(x+4|a|)=f(x)$。所以应写 $4|a|$，不要直接写 $4a$ 是正周期。

\medskip

\end{trapbox}
\clearpage
\begin{trapbox}
\textbf{\textcolor{CTrap}{坑五：忘记定义域}}
\begin{center}\begin{tikzpicture}[font=\small,>=stealth,line join=round]
\draw[gray!40,rounded corners=3pt] (0,0) rectangle (14.2,8.3);\node at (7.1,7.700000000000001) {只有非负输入时，负x可能根本没有函数值};
\begin{scope}[shift={(2.3857142857142857,3.8842857142857143)},x=0.8571428571428571cm,y=1.3142857142857143cm]\draw[->,gray] (-1.5,0)--(5.5,0) node[right]{$x$};\draw[->,gray] (0,-1.7)--(0,1.8) node[above]{$y$};\draw[blue,thick,solid,domain=0:1,samples=180] plot(\x,{1*\x+(0)});\draw[blue,thick,solid,domain=1:2,samples=180] plot(\x,{-1*\x+(2)});\draw[blue,thick,solid,domain=2:3,samples=180] plot(\x,{-1*\x+(2)});\draw[blue,thick,solid,domain=3:4,samples=180] plot(\x,{1*\x+(-4)});\fill[orange!85!black] (0,0) circle(2.2pt);\fill[orange!85!black] (4,0) circle(2.2pt);\end{scope}\node[anchor=west,align=left,text width=5.0cm] at (8.4,4.0) {$D=[0,4]$\\[9pt]图内有一段完整折线\\[9pt]但 $f(-1)$ 没有定义\\[9pt]无法对所有x代入 $-x$\\[9pt]也不能一直平移4};
\node[text width=12.6cm,align=center] at (7.1,.65) {这是定义域检查的反例；图像相似不等于函数已经具有奇偶性和周期。};
\end{tikzpicture}\end{center}



本结论使用时最方便的条件是函数定义在 $\mathbb{R}$ 上。比如只在 $x\ge0$ 上定义的函数，不可能在这个定义域上成为通常意义的奇函数或偶函数（除特殊定义方式外）；更不能对任意 $x$ 随便代入 $-x$。其他定义域要检查：$x\in D$ 时，$-x,2a-x$ 是否都在 $D$ 中；奇偶、对称等式是否对全部输入成立；向前和向后平移后定义域是否整体不变。只在图内少数点成立不能代替这些条件。

\medskip

\end{trapbox}
\clearpage
\begin{trapbox}
\textbf{\textcolor{CTrap}{坑六：直接把 $a=0$ 代入公式}}
\begin{center}\begin{tikzpicture}[font=\small,>=stealth,line join=round]
\draw[gray!40,rounded corners=3pt] (0,0) rectangle (14.2,8.3);\node at (7.1,7.700000000000001) {a等于0时，两种不同要求会迫使整条图像为0};
\begin{scope}[shift={(3.242857142857143,3.814705882352941)},x=0.8571428571428571cm,y=1.352941176470588cm]\draw[->,gray] (-2.5,0)--(4.5,0) node[right]{$x$};\draw[->,gray] (0,-1.6)--(0,1.8) node[above]{$y$};\draw[blue,thick,solid,domain=-2:4,samples=180] plot(\x,{0});\fill[orange!85!black] (0,0) circle(2.2pt);\fill[orange!85!black] (2,0) circle(2.2pt);\end{scope}\node[anchor=west,align=left,text width=5.0cm] at (8.4,4.0) {$f(-x)=f(x)$\\[9pt]同时 $f(-x)=-f(x)$\\[9pt]因此 $f\equiv0$\\[9pt]任意正数都是周期\\[9pt]0不是正周期};
\node[text width=12.6cm,align=center] at (7.1,.65) {零函数没有最小正周期；不能把4|a|=0当作公式给出的正周期。};
\end{tikzpicture}\end{center}



如果奇函数还关于 $x=0$ 对称，则 $f(-x)$ 既等于 $-f(x)$ 又等于 $f(x)$，故 $f(x)=0$；偶函数还关于 $(0,0)$ 对称也一样。零函数处处为零，每个正数都是它的周期，\textbf{不存在最小正周期}，而 $4|a|=0$ 不是正周期。相反，偶配轴或奇配点在 $a=0$ 时，新增条件可能完全没有增加信息；下文另给反例。

\medskip

\end{trapbox}
\clearpage
\begin{trapbox}
\textbf{\textcolor{CTrap}{坑七：两种搭配的“固定零点”记反}}
\begin{center}\begin{tikzpicture}[font=\small,>=stealth,line join=round]
\draw[gray!40,rounded corners=3pt] (0,0) rectangle (14.2,8.8);\node at (7.1,8.200000000000001) {先看哪一种对称强制取0，再沿反号关系继续找};
\node at (3.1500000000000004,7.3) {奇配轴：$f=\sin\frac{\pi x}2$};\begin{scope}[shift={(1.293939393939394,4.338709677419355)},x=0.7424242424242425cm,y=1.225806451612903cm]\draw[->,gray] (-0.8,0)--(5.8,0) node[right]{$x$};\draw[->,gray] (0,-1.5)--(0,1.6) node[above]{$y$};\draw[blue,thick,solid,domain=-0.5:5.5,samples=180] plot(\x,{sin(deg(pi*\x/2))});\fill[orange!85!black] (0,0) circle(2pt);\fill[orange!85!black] (2,0) circle(2pt);\fill[orange!85!black] (4,0) circle(2pt);\end{scope}\node[align=center,text width=5.8cm] at (3.1500000000000004,1.7) {奇函数先给 $f(0)=0$\\再得到 $f(2)=f(4)=0$};\node at (9.95,7.3) {偶配点：$g=\cos\frac{\pi x}2$};\begin{scope}[shift={(8.093939393939394,4.338709677419355)},x=0.7424242424242425cm,y=1.225806451612903cm]\draw[->,gray] (-0.8,0)--(5.8,0) node[right]{$x$};\draw[->,gray] (0,-1.5)--(0,1.6) node[above]{$y$};\draw[teal,thick,solid,domain=-0.5:5.5,samples=180] plot(\x,{cos(deg(pi*\x/2))});\fill[orange!85!black] (1,0) circle(2pt);\fill[orange!85!black] (3,0) circle(2pt);\fill[orange!85!black] (5,0) circle(2pt);\end{scope}\node[align=center,text width=5.8cm] at (9.95,1.7) {中心先给 $g(1)=0$\\再得到 $g(3)=g(5)=0$};
\node[text width=12.6cm,align=center] at (7.1,.65) {两图同尺度；奇配轴不能直接断言轴上的函数值为0，例如f(1)=1。};
\end{tikzpicture}\end{center}



奇函数先告诉你 $f(0)=0$，再通过 $f(x+2a)=-f(x)$ 得到 $f(2a)=0$、$f(4a)=0$ 等；偶函数\textbf{加点 $(a,0)$ 对称}先告诉你 $f(a)=0$，随后可得 $f(3a)=0$、$f(5a)=0$ 等。不能仅因“奇函数关于 $x=a$ 对称”就说 $f(a)=0$，例如 $\sin\frac{\pi x}{2}$ 在 $a=1$ 时取值为 $1$。

\medskip
\textbf{\textcolor{CTrap}{考场检查的三个小问题}}
\begin{enumerate}[leftmargin=*,itemsep=0.15em]
 \item 奇性或偶性，配的是\textbf{轴}还是\textbf{中心}？四种组合要分别识别。
 \item 推出来的是\textbf{变号}还是\textbf{不变号}？
 \item 题目问的是\textbf{一个周期、最小正周期，还是函数值}？
\end{enumerate}
\end{trapbox}
\clearpage
\begin{trapbox}
\textbf{坑八｜把轴 $x=a$ 和中心 $(a,0)$ 当成同一回事}
\begin{center}\begin{tikzpicture}[font=\small,>=stealth,line join=round]
\draw[gray!40,rounded corners=3pt] (0,0) rectangle (14.2,7.4);
\node[text width=12.6cm,align=center] at (7.1,6.75) {轴与中心不同：在同一条偶函数曲线上比较};

\begin{scope}[shift={(3.0,3.5)},x=1.9cm,y=1.3cm]
\draw[->,gray] (-.8,0)--(4.5,0) node[right]{$x$};
\draw[->,gray] (0,-1.35)--(0,1.55) node[above]{$y$};
\draw[blue,thick,domain=-.5:4.2,samples=180] plot(\x,{cos(deg(pi*\x/2))});
\draw[teal,dashed] (0,-1.2)--(0,1.2);\draw[teal,dashed] (2,-1.2)--(2,1.2);
\fill[violet] (1,0) circle(2.8pt);\node[above right=5pt,violet] at (1,0){中心 $(1,0)$};
\foreach \u/\v in {.5/.7071067812,1.5/-.7071067812}{\fill[orange!85!black] (\u,\v) circle(2.3pt);}
\node[above,teal] at (2,1.3){轴 $x=2$};
\end{scope}
\node at (7.1,1.35) {$f(\frac12)=\frac{\sqrt2}{2},\quad f(\frac32)=-\frac{\sqrt2}{2}$：两侧不是等高。};

\node[text width=12.6cm,align=center] at (7.1,.65) {$f(x)=\cos(\pi x/2)$；中心 $(1,0)$ 不能读成轴 $x=1$。};
\end{tikzpicture}\end{center}


\textbf{错误：}看到偶函数关于 $(1,0)$ 对称，就当作关于 $x=1$ 对称，写出 $2$ 是周期。
\textbf{反例：}图中两侧 $\tfrac12,\tfrac32$ 的高度相反而不相等，故 $x=1$ 不是轴；$f(0)=1$、$f(2)=-1$，$2$ 也不是周期。
\textbf{正确替代：}轴用 $f(2a-x)=f(x)$；零高度中心用 $f(2a-x)=-f(x)$。先识别它们，再选 $2|a|$ 或 $4|a|$。

\end{trapbox}
\clearpage
\begin{trapbox}
\textbf{坑九｜对称中心纵坐标不是零，也擅自删掉}
\begin{center}\begin{tikzpicture}[font=\small,>=stealth,line join=round]
\draw[gray!40,rounded corners=3pt] (0,0) rectangle (14.2,7.4);
\node[text width=12.6cm,align=center] at (7.1,6.75) {中心的高度不能丢：偶函数配 $(1,2)$ 的例子};

\begin{scope}[shift={(2.4,2.5)},x=1.55cm,y=.85cm]
\draw[->,gray] (-.8,0)--(5.3,0) node[right]{$x$};
\draw[->,gray] (0,-.3)--(0,3.7) node[above]{$y$};
\draw[gray,dashed] (-.5,2)--(5,2) node[right]{$y=2$};
\draw[blue,thick,domain=-.5:5,samples=200] plot(\x,{2+cos(deg(pi*\x/2))});
\fill[violet] (1,2) circle(2.7pt);\node[above right=4pt] at (1,2){$(1,2)$};
\foreach \u/\v in {0/3,2/1,4/3}{\fill[orange!85!black] (\u,\v) circle(2.3pt);\node[below=4pt] at (\u,0){$\u$};}
\end{scope}
\node at (7.1,1.35) {$f(0)=3,\quad f(2)=1\ne-3,\quad f(x+2)=4-f(x)$};

\node[text width=12.6cm,align=center] at (7.1,.65) {$f(x)=2+\cos(\pi x/2)$ 为偶函数；点对称使对应高度之和为 $4$，不是 $0$。};
\end{tikzpicture}\end{center}


\textbf{错误：}偶函数关于 $(a,c)$ 对称，就直接写 $f(x+2a)=-f(x)$。
\textbf{原因：}真实关系是 $f(2a-u)=2c-f(u)$，令 $u=-x$，由偶性得到 $f(x+2a)=2c-f(x)$。
图中 $c=2$，$f(2)=1$，并不等于 $-f(0)=-3$。
\textbf{正确替代：}F027 的点对称格要求中心是 $(a,0)$。若 $c\ne0$，须保留 $2c$；偶函数这一格再移 $2a$ 仍可还原，但中间关系不同。奇函数配非零高度中心不能照搬本表；一般双对称情况见 F026。

\end{trapbox}
\clearpage
\begin{trapbox}
\textbf{坑十｜$a=0$ 时，同类对称可能没有新增信息}
\begin{center}\begin{tikzpicture}[font=\small,>=stealth,line join=round]
\draw[gray!40,rounded corners=3pt] (0,0) rectangle (14.2,7.4);
\node[text width=12.6cm,align=center] at (7.1,6.75) {$a=0$：同类对称只是重复原条件};

\begin{scope}[shift={(3.0,2.5)},x=1.25cm,y=.7cm]
\draw[->,gray] (-2.3,0)--(2.5,0) node[right]{$x$};
\draw[->,gray] (0,-.3)--(0,4.8) node[above]{$y$};
\draw[blue,thick,domain=-2:2,samples=100] plot(\x,{\x*\x});
\draw[violet,dashed] (0,0)--(0,4.2);
\fill[orange!85!black] (0,0) circle(2.3pt);\fill[orange!85!black] (1,1) circle(2.3pt);
\end{scope}
\node[text width=5.6cm,align=left] at (10,3.9) {$f(x)=x^2$ 是偶函数，\\也关于 $x=0$ 对称。\\[8pt]若 $T>0$ 是周期，\\则 $f(T)=f(0)$，\\但 $T^2>0$，矛盾。};

\node[text width=12.6cm,align=center] at (7.1,.65) {没有新增不同的对称位置，就不能从 $2|a|=0$ 得到非零周期。};
\end{tikzpicture}\end{center}


\textbf{错误：}偶函数再关于 $x=0$ 对称，就把 $2|a|=0$ 当成周期，或认为它必有非零周期。
\textbf{反例：}$f(x)=x^2$ 本来就关于 $y$ 轴对称，重复这个条件不会产生周期。奇函数 $f(x)=x$ 关于 $(0,0)$ 对称也是同样道理；若 $T>0$ 是周期，则 $f(T)=f(0)$ 给出 $T=0$，矛盾。
\textbf{正确替代：}先核对 $a\ne0$。$a=0$ 时要分清：同类搭配可能只是原条件；不同类搭配会迫使函数恒为零，均不能把零当成正周期。

\end{trapbox}
